\chapter{Grundlagen}
\label{ch:grundlagen}
In diesem Kapitel werden die Grundlagen, die für das Verständnis der Arbeit benötigt werden, vermittelt. Im Folgenden werden zentrale Konzepte vorgestellt, auf denen die weiteren Kapitel aufbauen. Ziel dieses Kapitels ist es, ein allgemeines Verständnis aufzubauen, um eine fundierte Wissensbasis für die darauffolgenden Teile der Arbeit zu schaffen.

\section{Elektroenzephalographie (EEG)}
\label{sec:eeg}

Die Elektroenzephalographie (EEG) ist ein Verfahren zur Aufzeichnung der elektrischen Aktivität des menschlichen Gehirns. Die aufgezeichneten Signale bestehen aus Spannungsschwankungen im Mikrovoltbereich, welche aus der synchronen elektrischen Aktivität von Neuronen resultieren \autocite{blum2017}. Die Messung erfolgt über Elektroden, die auf der Kopfhaut platziert werden und die Potenzialdifferenzen zwischen verschiedenen Punkten erfassen \autocite{10353995}.

\paragraph{Eigenschaften des Signals}
EEG-Signale besitzen eine äußerst geringe Amplitude, die in der Regel im Bereich von 2 bis 100 Mikrovolt (µV) liegt. Das Frequenzspektrum erstreckt sich dabei meist von 0,1~Hz bis 100~Hz \autocite{zhao2025eeg,10353995}.

\paragraph{Herausforderungen durch Rauschen und Artefakte}
In seiner Rohform ist ein EEG-Signal stark von Rauschen geprägt, weshalb es ein unzureichendes Signal-Rausch-Verhältnis (SNR) aufweist \autocite{blum2017}.
Da die erwarteten neuronalen Aktivitäten schwach sind, werden sie schnell von amplitudenstarken physiologischen Artefakten (wie z.\,B. Muskel- oder Augenbewegungen) sowie technischen Störeinflüssen überlagert, was zu Verzerrungen des Originalsignals führt \autocite{10353995}.
% TODO: Was ist EEG, Messung, Signalform (Kanäle, µV, 250 Hz)  \autocite{zhao2025eeg}

\section{Frequenzbänder}
\label{sec:baender}
% TODO: Delta–Gamma, Bedeutung, Bandgrenzen, Variabilität in der Literatur  \autocite{sanei2013}

\section{Frequenzanalyse}
\label{sec:frequenzanalyse}
% TODO: FFT, Spektrum/Band-Power, spektrale Leckage + Hann \autocite{harris1978}, Welch \autocite{welch1967}

\section{Lokale Sprachmodelle auf mobilen Geräten}
\label{sec:llm-grundlagen}
% TODO: was sind LLMs, On-Device-Inferenz, Quantisierung  \autocite{smalllm2024,ondevice2024}